\input{../tex-inputs/latex-leseansicht-vorspann}

\section[Arthur Schnitzler an Gustav Schwarzkopf, 28. 3. 1903]{L04062 Arthur Schnitzler an Gustav Schwarzkopf, 28. 3. 1903}
\nopagebreak\mylabel{L04062v}
\rehead{ }\normalsize\beginnumbering\briefempfaengerindex{Schwarzkopf, Gustav@\textsc{Schwarzkopf, Gustav}!zzzSchnitzler, Arthur@\emph{von Arthur Schnitzler}!1903-03-282@{28. 3. 1903}|(be}
\toendnotes[C]{\smallbreak\pagebreak[2]}
\correspDesc{Versand  durch Arthur Schnitzler am 28. 3. 1903 in Wien
\newline{}Erhalt  durch Gustav Schwarzkopf im Zeitraum [28. 3. 1903 – 31. 3. 1903?] in Wien}\toendnotes[C]{\smallbreak}
\Standort{CUL, Schnitzler, B 96.}
\physDesc{Brief, 1 Blatt, 1 Seite, 173 Zeichen
\newline{}Handschrift: schwarze Tinte, deutsche Kurrent}\toendnotes[C]{\smallbreak}
\pstart
           \noindent{}{\pb}lieber Guſtav, ich bitte
      Sie herzlich, am \label{K_L04062-1v}\edtext{Montag
               mit uns in den Circus\oindex{Wien@\textbf{Wien}!II., Leopoldstadt@\textbf{II., Leopoldstadt}!Zirkus Renz@\textbf{Zirkus Renz}, \emph{Veranstaltungsgebäude}|pwv}\eventindex{Zirkus Renz@\textbf{Zirkus Renz}!Zirkus Schumann, 30.3.1903@Zirkus Schumann, 30.3.1903|pwv}}{\lemma{\textnormal{\emph{Montag … Circus}}}\Cendnote{\textnormal{Vgl. A. S.: \emph{Kulturveranstaltungen}, 30. 3. 1903.}}}\label{K_L04062-1} zu
      gehen, beifolgend das Billet
      und viele Grüße.\pend
           
\pstart
           \label{K_L04062-2v}\edtext{Von
               W.s\pwindex{Wassermann, Jakob 10.\,3.\,1873 Fürth – 1.\,1.\,1934 Altaussee@\textsc{Wassermann, Jakob} (10.\,3.\,1873 Fürth – 1.\,1.\,1934 Altaussee), \emph{Schriftsteller}|pwu}\pwindex{Wassermann, Julie 5.\,12.\,1876 Wien – April 1963 Zürich@\textsc{Wassermann, Julie} (5.\,12.\,1876 Wien – April 1963 Zürich), \emph{Schriftstellerin}|pwu} ſind Sie durch Olga\pwindex{Schnitzler, Olga 17.\,1.\,1882 Wien – 13.\,1.\,1970 Lugano@\textsc{Schnitzler, Olga} (17.\,1.\,1882 Wien – 13.\,1.\,1970 Lugano), \emph{Schauspielerin, Sängerin}|pw} u
               mich geſchieden.}{\lemma{\textnormal{\emph{Von … geschieden.}}}\Cendnote{\textnormal{In Schnitzlers{ }\emph{Tagebuch}\pwindex{Schnitzler, Arthur 15.\,5.\,1862 Wien – 21.\,10.\,1931 ebd.@\textsc{Schnitzler, Arthur} (15.\,5.\,1862 Wien – 21.\,10.\,1931 ebd.), \emph{Schriftsteller, Mediziner}!Tagebuch@\strich\emph{Tagebuch}|pwk}
                  finden sich – neben Schwarzkopf\pwindex{Schwarzkopf, Gustav 7.\,11.\,1853 Wien – 13.\,11.\,1939 ebd.@\textsc{Schwarzkopf, Gustav} (7.\,11.\,1853 Wien – 13.\,11.\,1939 ebd.), \emph{Schriftsteller}|pwk} – an diesem Tag keine weiteren Besucher beim Zirkus Schumann\eventindex{Zirkus Renz@\textbf{Zirkus Renz}!Zirkus Schumann, 30.3.1903@Zirkus Schumann, 30.3.1903|pwk} namentlich genannt.
                  Zu einer Animosität zwischen Schwarzkopf\pwindex{Schwarzkopf, Gustav 7.\,11.\,1853 Wien – 13.\,11.\,1939 ebd.@\textsc{Schwarzkopf, Gustav} (7.\,11.\,1853 Wien – 13.\,11.\,1939 ebd.), \emph{Schriftsteller}|pwk} und Jakob\pwindex{Wassermann, Jakob 10.\,3.\,1873 Fürth – 1.\,1.\,1934 Altaussee@\textsc{Wassermann, Jakob} (10.\,3.\,1873 Fürth – 1.\,1.\,1934 Altaussee), \emph{Schriftsteller}|pwk} und 
                  Julie Wassermann\pwindex{Wassermann, Julie 5.\,12.\,1876 Wien – April 1963 Zürich@\textsc{Wassermann, Julie} (5.\,12.\,1876 Wien – April 1963 Zürich), \emph{Schriftstellerin}|pwk}, von denen die Rede sein dürfte, gibt es keine Überlieferung.}}}\label{K_L04062-2}\pend
           
\pstart
           Ihr{\\[\baselineskip]}\spacefill\mbox{A.}\pend
           \leftskip=0em{}
\pstart
           28. 3. 903.\pend
           \selectlanguage{ngerman}\endnumbering\briefempfaengerindex{Schwarzkopf, Gustav@\textsc{Schwarzkopf, Gustav}!zzzSchnitzler, Arthur@\emph{von Arthur Schnitzler}!1903-03-282@{28. 3. 1903}|)be}\mylabel{L04062h}
\begin{anhang}
\end{anhang}\newcommand{\dateiname}{L04062}\newcommand{\titel}{Arthur Schnitzler an Gustav Schwarzkopf, 28. 3. 1903}\newcommand{\editorInnen}{Herausgegeben von Jahnke, SelmaMüller, Martin Anton}\input{../tex-inputs/latex-leseansicht-abspann}
